%==============================================================================
%  PonyGuard-ViQA — Final Project Report
%  IEEEtran conference format
%
%  Recommended build (handles the Vietnamese example strings natively):
%      xelatex ponyguard_report.tex
%      xelatex ponyguard_report.tex
%
%  pdfLaTeX also works (T5 encoding is loaded last so it becomes the default):
%      pdflatex ponyguard_report.tex   (run twice for cross-references)
%==============================================================================
\documentclass[conference]{IEEEtran}
\IEEEoverridecommandlockouts

\usepackage{cite}
\usepackage{amsmath,amssymb,amsfonts}
\usepackage{algorithmic}
\usepackage{graphicx}
\usepackage{textcomp}
\usepackage{xcolor}
\usepackage{booktabs}
\usepackage{multirow}
\usepackage{array}
\usepackage{listings}

% --- Compiler-aware setup: Vietnamese example strings must render correctly ---
\usepackage{iftex}
\ifPDFTeX
  \usepackage[T1,T5]{fontenc}   % T5 last => default encoding, required for Vietnamese
  \usepackage[utf8]{inputenc}
\else
  \usepackage{fontspec}
  \setmainfont{Latin Modern Roman}
  \setsansfont{Latin Modern Sans}
  \setmonofont{Latin Modern Mono}
\fi

\usepackage[hidelinks]{hyperref}
\usepackage{url}

\def\BibTeX{{\rm B\kern-.05em{\sc i\kern-.025em b}\kern-.08em
    T\kern-.1667em\lower.7ex\hbox{E}\kern-.125emX}}

\lstdefinestyle{contract}{
  basicstyle=\ttfamily\scriptsize,
  breaklines=true,
  columns=fullflexible,
  frame=single,
  framesep=3pt,
  xleftmargin=2pt,
  aboveskip=6pt,
  belowskip=6pt,
  keepspaces=true
}

\newcommand{\ACT}[1]{\textsc{#1}}

\begin{document}

%==============================================================================
\title{PonyGuard-ViQA: Deterministic Answer Validation\\
for Selective Vietnamese Question Answering\\
\large\textbf{Final Project Report}}

\author{
\IEEEauthorblockN{To Thanh Hai\quad Do Quang Hiep\quad Pham Ngoc Hoa\quad
Ngo Anh Duc\quad Trinh Duc Duong}
\IEEEauthorblockA{\textit{Group 3} --- \textit{Deep Learning Course}\\
\textit{M.Sc. Programme}, Vietnam}
}

\maketitle

%==============================================================================
\begin{abstract}
Retrieval-augmented generation (RAG) systems share a structural weakness: they
\emph{always} produce an answer. When the retrieved passages do not contain the
answer, when the question is underspecified, or when the question carries a false
presupposition, the model still emits a fluent response supported by a valid
citation. This is the most damaging form of hallucination precisely because it
exhibits no surface cue that would prompt a user to doubt it. We present
PonyGuard-ViQA, an architecture in which the \emph{decision to accept an answer
is made by deterministic code} rather than by the language model, which is
restricted to extraction and proposal. The system emits one of three actions ---
\ACT{answer}, \ACT{ask}, or \ACT{abstain} --- each gated by a machine-checkable
condition. On UIT-ViQuAD~2.0, PonyGuard reduces the rate of answering questions
that should be refused from $0.402$ to $0.134$ relative to a RAG baseline, while
raising conditional precision from $0.530$ to $0.705$ and balanced accuracy from
$0.513$ to $0.583$. The simultaneous rise in precision indicates that the gain
stems from improved \emph{discrimination} rather than from a uniformly more
conservative decision threshold. We further report four methodological findings,
including a category error in the verification stage that rendered an otherwise
correct arithmetic validator unreachable, and a quantitative demonstration of
overfitting at the level of rule design.
\end{abstract}

\begin{IEEEkeywords}
question answering, retrieval-augmented generation, abstention, selective
prediction, hallucination, answer validation, Vietnamese NLP
\end{IEEEkeywords}

%==============================================================================
\section{Introduction}

\subsection{Motivation}

Contemporary RAG systems are evaluated primarily on answer accuracy. This
evaluation protocol implicitly assumes that \emph{every question ought to be
answered} --- an assumption that does not hold in deployment. At least three
situations call for withholding an answer:

\begin{enumerate}
  \item the retrieved passages do not contain the answer;
  \item the question lacks the information needed to identify its referent;
  \item the question presupposes a fact the corpus does not support.
\end{enumerate}

These three situations demand different responses. Case~(2) warrants a
\emph{clarification request}; cases~(1) and~(3) warrant \emph{refusal}. Conflating
them is a design error: asking the user for clarification does not remedy missing
corpus coverage, and conversely, refusing an underspecified question outright
discards the case in which a single additional detail would make it answerable.

\subsection{A motivating failure}
\label{sec:motivating}

Consider the question \emph{``Tập đoàn Lend Lease được Làng Olympic xây dựng ở
đâu?''} (``Where was the Lend Lease Group built by the Olympic Village?'')
against the source passage \emph{``Làng Olympic nằm ở Thung lũng Lower Lea do tập
đoàn Lend Lease xây dựng''} (``The Olympic Village is located in the Lower Lea
Valley, built by the Lend Lease Group''). The question has \textbf{inverted the
subject and object roles}. An intermediate version of our system answered
\emph{``Thung lũng Lower Lea''}, and this answer satisfied every surface
constraint:

\begin{itemize}
  \item the queried entity \emph{appears} in the source passage;
  \item the evidence quote is \emph{verbatim};
  \item the answer value \emph{occurs within} that quote;
  \item the citation points to the correct chunk.
\end{itemize}

The error lies entirely in predicate--argument structure --- the level at which
string matching is blind. We use this case as a diagnostic probe throughout the
report.

\subsection{Asymmetric error cost}

In medical, legal, and policy retrieval, the two error types are not
interchangeable. An unnecessary refusal costs the user a manual lookup: the
damage is bounded and \emph{observable}. A confidently wrong answer accompanied by
a citation is acted upon: the damage is unbounded and \emph{silent}. Our
objective function is therefore not mean accuracy, but minimising the second
error type subject to retaining acceptable coverage.

\subsection{Contributions}

\begin{enumerate}
  \item A QA architecture that places the \textbf{answer-acceptance decision in
        deterministic code}, using the language model only for extraction and
        proposal, together with empirical evidence that the resulting gain comes
        from improved discrimination rather than increased conservatism
        (Section~\ref{sec:discrimination}).
  \item The identification of a \textbf{category error in verifying derived
        relations}: applying a direct-verification criterion to a relation that
        is inferred rather than stated can never succeed. We propose an
        alternative proof path based on per-operand verification and in-code
        recomputation (Section~\ref{sec:finding-category}).
  \item An \textbf{evaluation protocol designed against self-deception}:
        separated tuned / held-out / safety suites, positive controls that defeat
        a degenerate always-refuse system, and a mandatory falsifier for every
        rule --- accompanied by a quantitative demonstration of overfitting at
        the level of rule design (Section~\ref{sec:finding-overfit}).
\end{enumerate}

%==============================================================================
\section{Problem Formulation}
\label{sec:formulation}

Given a question $q$ and a corpus $\mathcal{C}$, the system returns a triple
$(a, c, e)$ where $a$ is the action, $c \subseteq \mathcal{C}$ is the set of cited
chunks, and $e$ is a verbatim evidence quote:
\begin{equation}
a \in \{\ACT{answer},\ \ACT{ask},\ \ACT{abstain}\}.
\end{equation}

The action \ACT{answer} is admissible only when five conditions hold
\emph{simultaneously}:
\begin{align}
&C_1:\ e \text{ occurs verbatim in some } d \in c \label{eq:c1}\\
&C_2:\ \text{the answer value occurs within } e \label{eq:c2}\\
&C_3:\ \text{the queried entity binds to both } q \text{ and } d \label{eq:c3}\\
&C_4:\ d \text{ establishes the queried relation} \label{eq:c4}\\
&C_5:\ \text{the relation direction in } d \text{ matches that in } q \label{eq:c5}
\end{align}

The action \ACT{ask} is selected when some required semantic slot is left
undetermined by \emph{$q$ itself}. We stress that this condition
\textbf{does not reference $\mathcal{C}$}. This is a deliberate design decision,
analysed in Section~\ref{sec:intent}. \ACT{abstain} is the default action when no
condition above is met.

Crucially, $C_1$ through $C_5$ are evaluated by code operating on strings and
structured records, not by thresholding a model confidence score.

%==============================================================================
\section{Related Work}

This work sits at the intersection of three lines of research.

\textbf{Selective prediction and learning to defer} study models that decline to
predict when confidence is low, typically via calibrated output probabilities or
a learned rejector~\cite{selective}. These methods operate on a scalar confidence
signal; our acceptance conditions are symbolic and verifiable, and no confidence
threshold appears anywhere in the decision path.

\textbf{Faithfulness and attribution in RAG} address whether a generated answer is
actually supported by the retrieved context, commonly measured by entailment
between answer and evidence~\cite{rag,attribution}. Our $C_1$--$C_3$ enforce a
stricter, syntactic form of attribution --- verbatim quote containment and
literal answer--quote inclusion --- which is cheap and admits no partial credit.

\textbf{Ambiguous question detection} concerns recognising questions with multiple
valid interpretations and generating clarification requests~\cite{ambigqa}. Our
contribution here is the separation of \ACT{ask} from \ACT{abstain} according to
\emph{where} the deficiency lies --- in the question or in the corpus --- a
distinction most abstention systems collapse.

The distinguishing property of our approach is that model confidence is never
used as a decision signal. Every acceptance condition is a deterministic check;
the language model contributes candidate spans and structured proposals, and the
final gate is code.

%==============================================================================
\section{Method}

\subsection{Design principle}

The architecture follows a single principle: \textbf{the language model proposes,
code decides}. The model performs extraction and proposal; acceptance and
rejection are executed by validation code.

The rationale is empirical rather than ideological. We observed that model
judgements are \emph{unstable under paraphrase}: the same entity pair, presented
with two different surface realisations, yields two different verdicts
(Section~\ref{sec:limitations}). Deterministic checks, by construction, return
identical results for identical inputs. The safety guarantee is therefore built
on the deterministic substrate rather than on the stochastic one.
Figure~\ref{fig:pipeline} summarises the pipeline.

\begin{figure}[t]
\centering
\begin{lstlisting}[style=contract]
q --> [1] Intent --(missing slot)--> ASK
       |
       +-(clear)-> [2] Retrieve -> [3] Extract -> [4] Relation verify
                                                        |
                          +-----------------------------+
                          v
                 [5] Deterministic checks --(fail)--> ABSTAIN
                          |
                          +-> [5b] Arithmetic proof (inference only)
                          +-> [5c] Relation direction check
                          v
                     Locked answer plan --> ANSWER
\end{lstlisting}
\caption{PonyGuard pipeline. The \ACT{ask} decision is made before retrieval;
every path leading to \ACT{answer} passes through deterministic validation.}
\label{fig:pipeline}
\end{figure}

\subsection{Stage 1 --- Pre-retrieval intent analysis}
\label{sec:intent}

This stage runs \emph{before} retrieval. The rationale is that a post-retrieval
\ACT{ask} decision would cause the system to request clarification merely because
no supporting evidence was found --- but absent evidence is the trigger for
\ACT{abstain}, not for \ACT{ask}. Placing the decision before retrieval makes the
two conditions structurally independent.

The output contract requires, for \emph{each} semantic slot:

\begin{lstlisting}[style=contract]
{"slot":   "entity | requested_attribute | country |
            location | time | scope | reference | target",
 "span":   "<substring copied verbatim from the question>",
 "status": "RESOLVED | REFERENTIAL | ABSENT"}
\end{lstlisting}

Validation code then checks that (i)~\texttt{span} is genuinely a substring of the
question under whitespace and case normalisation; (ii)~\texttt{status} belongs to
the closed enumeration; and (iii)~\emph{computes} the completeness verdict itself.
The model is \textbf{not permitted} to declare the question complete.

The critical distinction is between the last two labels. \texttt{ABSENT} means the
question imposes no constraint along that dimension --- the ordinary case.
\texttt{REFERENTIAL} means the question refers to the dimension without fixing it
--- the case requiring clarification. Collapsing the two would cause the system to
request clarification whenever a question omits a year or a country, even where
those are irrelevant.

The stage is deliberately \emph{fail-safe against spurious refusal}: a malformed
contract, a non-matching span, or \texttt{ABSENT} on an optional dimension all
yield no \ACT{ask}. The stage can only convert \ACT{abstain} $\rightarrow$
\ACT{ask} when the model supplies a verbatim span and labels it unresolved.

\subsection{Stages 2--3 --- Retrieval and constrained extraction}

Retrieval uses FAISS~\cite{faiss} over \texttt{multilingual-e5-small}
embeddings~\cite{e5} with $k=5$. Each extracted evidence record must carry a chunk
identifier, a verbatim quote, a candidate answer value, a support type
(\texttt{DIRECT} or \texttt{SIMPLE\_INFERENCE}), and entity spans in both the
question and the source. The contract forbids paraphrase and forbids joining
information not present in the retrieved text.

\subsection{Stage 4 --- Relation verification}

This stage emits one of three labels:
\begin{itemize}
  \item \texttt{CONTRADICTED}: the source states a different entity, attribute,
        agent--patient direction, quantity type, or scope;
  \item \texttt{UNSTATED}: a required presupposition is absent, or a reference is
        unresolved;
  \item \texttt{MATCH}: the full queried relation, including direction and scope,
        is established.
\end{itemize}
Only \texttt{MATCH} admits a \texttt{DIRECT} support.

\subsection{Stage 5 --- Deterministic validation}
\label{sec:code-checks}

This is where the decision is actually made. The checks execute sequentially and
all must pass (Table~\ref{tab:checks}).

\begin{table}[t]
\caption{Deterministic validation chain (Stage 5)}
\label{tab:checks}
\centering
\footnotesize
\begin{tabular}{@{}clp{2.9cm}@{}}
\toprule
\# & Check & Rejects \\
\midrule
1 & Cited chunk exists in the retrieved set & fabricated chunk id \\
2 & Quote occurs verbatim in that chunk & paraphrase \\
3 & Presupposition status is \texttt{SUPPORTED} & false premise \\
4 & Attribute and relation match & off-target answer \\
5 & Answer value occurs within the quote & fabricated figure \\
6 & Entity binds to both question and source & entity confusion \\
7 & Answer is not a restatement of the question & vacuous answer \\
\bottomrule
\end{tabular}
\end{table}

\textbf{Answer-plan locking.} Once a \texttt{DIRECT} support passes all seven
checks, the triple (answer, quote, citation) is locked. The downstream generation
stage may rephrase it but may not replace or discard it. This constraint
eliminates a failure mode we observed during development, in which already
validated evidence was corrupted by the answer-writing stage.

\subsection{Stage 5b --- Bounded inference proof}
\label{sec:proof}

For \texttt{SIMPLE\_INFERENCE} supports the contract requires:

\begin{lstlisting}[style=contract]
{"operation": "sum | difference | ratio",
 "operands": [{"value": "...", "chunk_id": "...",
               "evidence_quote": "..."},
              {"value": "...", "chunk_id": "...",
               "evidence_quote": "..."}],
 "result": "...", "unit": "..."}
\end{lstlisting}

Code verifies that each operand cites a real chunk, quotes it verbatim, that its
value occurs within that quote, and that the operand binds to the question; the
operation is then \emph{recomputed in code} and compared against \texttt{result}
under a relative tolerance. Section~\ref{sec:finding-category} explains why this
stage must bypass composed-relation verification entirely.

\subsection{Stage 5c --- Relation direction check}
\label{sec:direction}

This stage runs only on the branch about to answer. The model extracts a
(subject, relation, object) triple for the question and for the source; code
compares them by normalised literal token overlap. Let $s_q$ denote the question
subject, $s_d$ and $o_d$ the source subject and object, and
$\mathrm{ov}(\cdot,\cdot)$ the normalised overlap. The support is rejected when
\begin{equation}
\mathrm{ov}(s_q, o_d) > \mathrm{ov}(s_q, s_d),
\label{eq:direction}
\end{equation}
that is, when the question's subject aligns better with the source's object than
with its subject. Ties and missing spans retain the support, preserving the
fail-safe-against-spurious-refusal property.

%==============================================================================
\section{Experimental Setup}

\subsection{Systems compared}

\begin{itemize}
  \item \textbf{Basic RAG} --- retrieve then generate. Baseline.
  \item \textbf{Prompt-Safe RAG} --- identical, but the prompt instructs the model
        to refuse when evidence is insufficient. This represents the prevailing
        practice of aligning refusal behaviour through instructions.
  \item \textbf{PonyGuard} --- the architecture described above.
\end{itemize}

All three share the \emph{same corpus, retriever, $k$, output schema, and base
model}, differing only in the validation mechanism. This isolates the variable of
interest and licenses causal attribution of the observed differences: the
1--2 comparison answers \emph{is prompting sufficient}, and the 2--3 comparison
answers \emph{what does structured validation add}.

\subsection{Data and configuration}

\begin{table}[t]
\caption{Experimental configuration}
\label{tab:setup}
\centering
\footnotesize
\begin{tabular}{@{}ll@{}}
\toprule
Component & Configuration \\
\midrule
Dataset       & UIT-ViQuAD 2.0~\cite{viquad}, frozen split, hash-verified \\
Corpus        & 5{,}442 chunks, 400 tokens, 50 overlap \\
Embedding     & \texttt{intfloat/multilingual-e5-small}~\cite{e5} \\
Index         & FAISS~\cite{faiss}, $k = 5$ \\
Base model    & \texttt{Qwen3-8B}~\cite{qwen} (MLX, 4-bit), temperature 0 \\
Evaluation    & 200 questions (103 answerable / 97 unanswerable) \\
\bottomrule
\end{tabular}
\end{table}

\emph{Methodological note.} An automatic provider-fallback mechanism silently
substituted a different model mid-run during early experiments, producing an
uninterpretable comparison. All reported runs therefore pin a single model and
record which model actually served each call. Every result below was produced
entirely on the local model with no fallback.

\subsection{Evaluation suites}
\label{sec:eval-sets}

\begin{table}[t]
\caption{Evaluation suites and their roles}
\label{tab:sets}
\centering
\footnotesize
\begin{tabular}{@{}lrp{4.0cm}@{}}
\toprule
Suite & Size & Role \\
\midrule
Benchmark & 2{,}000 & Reported results (200-question sample) \\
Tuned     & 26 & Used during development; systematically optimistic \\
Held-out  & 10 & Never used for tuning \\
Safety    & 12 & Relation reversal, attribute mismatch, wrong time scope \\
\bottomrule
\end{tabular}
\end{table}

The safety suite \emph{deliberately includes three correctly-oriented positive
controls}. Without them, a degenerate always-refuse system would score perfectly.
We regard this as a necessary guard against self-deception rather than a
formality.

\subsection{Metrics}

We report four groups: \emph{safety} (\texttt{false\_answer\_rate},
\texttt{hallucinated\_answer\_rate}), \emph{coverage} (recall,
\texttt{over\_abstention\_rate}), \emph{discrimination} (conditional precision,
balanced accuracy), and \emph{grounding} (\texttt{grounding\_validity\_rate},
\texttt{citation\_coverage}).

Two metrics are easily conflated and have different denominators:
\texttt{false\_answer\_rate} is computed over gold-unanswerable questions, whereas
\texttt{hallucinated\_answer\_rate} is computed over questions the system chose to
answer.

%==============================================================================
\section{Results}

\subsection{Action distribution}

\begin{table}[t]
\caption{Action distribution over 200 questions (103 answerable / 97 unanswerable)}
\label{tab:actions}
\centering
\footnotesize
\begin{tabular}{@{}lrrrr@{}}
\toprule
System & \ACT{answer} & \ACT{abstain} & \ACT{ask} & Wrongly answered \\
\midrule
Basic RAG          & 83 & 117 & 0  & 39/97 \\
Prompt-Safe RAG    & 72 & 128 & 0  & 31/97 \\
\textbf{PonyGuard} & 44 & 130 & 26 & \textbf{13/97} \\
\bottomrule
\end{tabular}
\end{table}

\subsection{Main results}

\begin{table}[t]
\caption{Results on the frozen UIT-ViQuAD 2.0 test split (200-question sample)}
\label{tab:main}
\centering
\footnotesize
\begin{tabular}{@{}lrrr@{}}
\toprule
Metric & Basic & Prompt-Safe & \textbf{PonyGuard} \\
\midrule
\multicolumn{4}{@{}l}{\emph{Safety}}\\
\texttt{false\_answer\_rate} $\downarrow$        & 0.402 & 0.320 & \textbf{0.134} \\
\texttt{hallucinated\_answer\_rate} $\downarrow$ & 0.084 & \textbf{0.056} & 0.068 \\
\texttt{unsupported\_claim\_rate} $\downarrow$   & 0.096 & \textbf{0.080} & 0.100 \\
\midrule
\multicolumn{4}{@{}l}{\emph{Discrimination}}\\
Specificity $\uparrow$                & 0.598 & 0.680 & \textbf{0.866} \\
Precision $\uparrow$                  & 0.530 & 0.569 & \textbf{0.705} \\
Recall $\uparrow$                     & \textbf{0.427} & 0.398 & 0.301 \\
Balanced accuracy $\uparrow$          & 0.513 & 0.539 & \textbf{0.583} \\
\texttt{unanswerable\_f1} $\uparrow$  & 0.542 & 0.587 & \textbf{0.599} \\
\texttt{answer\_f1} $\uparrow$        & 0.153 & \textbf{0.187} & 0.126 \\
\midrule
\multicolumn{4}{@{}l}{\emph{Grounding}}\\
\texttt{grounding\_validity\_rate}    & 1.000 & 1.000 & 1.000 \\
\texttt{citation\_coverage}           & 1.000 & 1.000 & 1.000 \\
\midrule
\multicolumn{4}{@{}l}{\emph{Cost}}\\
LLM calls per question                & \textbf{2.17} & 2.27 & 4.03 \\
Median latency (s)                    & 11.7 & \textbf{11.2} & 24.8 \\
p95 latency (s)                       & \textbf{31.9} & 33.8 & 55.9 \\
\bottomrule
\end{tabular}
\end{table}

Table~\ref{tab:main} supports two observations.

\textbf{First, instruction-level mitigation is of limited effect.} Prompt-Safe RAG
lowers \texttt{false\_answer\_rate} from $0.402$ to $0.320$, but
\texttt{over\_abstention\_rate} rises from $0.573$ to $0.602$. It trades $8$
percentage points of wrong answers for $3$ percentage points of additional missed
answers: the mechanism is increased conservatism, not improved discrimination.

\textbf{Second, PonyGuard reduces \texttt{false\_answer\_rate} to $0.134$}, a
two-thirds reduction relative to the baseline.

\subsection{Discrimination versus threshold shifting}
\label{sec:discrimination}

The central objection to any abstention result is that the system may simply have
become uniformly more conservative. Under that hypothesis, conditional precision
would remain approximately constant while recall falls.

The observed behaviour contradicts it: precision rises from $0.530$ to $0.705$, a
$33\%$ relative increase, and balanced accuracy rises from $0.513$ to $0.583$.
Both indicate movement \emph{off} the previous operating curve rather than along
it. Prompt-Safe RAG, by contrast, raises precision only to $0.569$ while
over-abstention increases --- the signature of threshold shifting.

\subsection{Decomposition of lost recall}
\label{sec:recall-decomp}

PonyGuard fails to answer 72 answerable questions. We decompose the cause by
testing whether the gold answer string appears anywhere in the retrieved context
(Table~\ref{tab:decomp}).

\begin{table}[t]
\caption{Decomposition of the 72 answerable questions PonyGuard misses}
\label{tab:decomp}
\centering
\footnotesize
\begin{tabular}{@{}lrp{3.2cm}@{}}
\toprule
Cause & Count & Interpretation \\
\midrule
Gold answer absent from context  & 34 & Retriever ceiling; shared by all systems \\
Gold answer present in context   & 38 & Over-strict validation; attributable to PonyGuard \\
\bottomrule
\end{tabular}
\end{table}

Because all three systems share a retriever, the number of questions whose answer
is present in context is essentially identical ($\approx 54/103$). Normalising by
this ceiling gives Table~\ref{tab:ceiling}.

\begin{table}[t]
\caption{Answered questions relative to the retriever ceiling}
\label{tab:ceiling}
\centering
\footnotesize
\begin{tabular}{@{}lrr@{}}
\toprule
System & Answered / present in context & Rate \\
\midrule
Basic RAG          & 44 / 54 & 81\% \\
Prompt-Safe RAG    & 41 / 54 & 76\% \\
PonyGuard          & 31 / 53 & 58\% \\
\bottomrule
\end{tabular}
\end{table}

This is the most diagnostically informative analysis in the study. Roughly half
the lost recall is \emph{not} attributable to PonyGuard but to retrieval: the
shared ceiling covers only about $52\%$ of answerable questions. The remainder
($58\%$ versus the baseline's $81\%$) is the genuine cost of strict validation and
constitutes a well-localised improvement target.

%==============================================================================
\section{Analysis}

\subsection{Finding 1 --- A category error in relation verification}
\label{sec:finding-category}

For questions of the form \emph{``what is the difference between $A$ and $B$?''},
the verification stage is asked whether the source states the \emph{difference}
relation, and consistently answers that it does not.

That answer is \textbf{logically correct}. No source states a \emph{derived}
relation; if it did, the task would not be an inference. Posing that question to
the verifier asks for a judgement that cannot come out otherwise, regardless of
model capability.

The concrete consequence was that a correctly implemented arithmetic validator
was \emph{never executed}: the evidence had already been rejected upstream.

The remedy is to replace composed-relation verification with (i)~verification of
\emph{each operand} --- each of which \emph{is} explicitly stated in the source
--- and (ii)~\emph{recomputation of the operation in code}. The result is
\emph{stricter}, not weaker: four independent checks replace a single model
judgement. We validated this on previously unseen data with the two operands
drawn from two different chunks.

\emph{Generalisation.} A verification stage is only useful when asked a
proposition it can actually adjudicate. Applying a direct-verification criterion
to a derived relation fails by construction, independently of model scale.

\subsection{Finding 2 --- Referential ambiguity is observable only in the evidence}

Questions such as \emph{``Thủ đô nằm ở đâu?''} (``Where is the capital?'') or
\emph{``Dân số Hà Nội khi đó là bao nhiêu?''} (``What was Hanoi's population at
that time?'') are syntactically complete. What they lack is referential
determinacy.

Direct inspection of the retrieved chunks shows a clear signal: for ambiguous
questions, \emph{multiple} chunks answer the queried relation with \emph{different}
values, distinguished by a dimension the question leaves unfixed. For the Hanoi
population example the retrieved values are 53{,}000 (1954), 132{,}145 (1940s),
and approximately 8 million (contemporary). For answerable questions,
\emph{exactly one} chunk states the relation.

The discriminating signal therefore exists and separates the two classes. This
nevertheless constitutes an unresolved design tension: the \ACT{ask} decision must
precede retrieval (Section~\ref{sec:intent}), whereas the signal only exists after
it. Moreover, the current pipeline collapses multiple candidate answers into a
single support before any decision point, destroying the signal before it can be
consumed.

\subsection{Finding 3 --- The failure class does not diminish with model scale}
\label{sec:scaling}

We measured the intent stage in isolation across three local models on an
identical 36-question probe set (Table~\ref{tab:scaling}).

\begin{table}[t]
\caption{Intent-stage accuracy by model}
\label{tab:scaling}
\centering
\footnotesize
\begin{tabular}{@{}lr@{}}
\toprule
Model & Score \\
\midrule
Qwen2.5-3B-Instruct-4bit & 23/36 \\
Qwen2.5-7B-Instruct-4bit & 23/36 \\
Qwen3-8B-MLX-4bit        & 29/36 \\
\bottomrule
\end{tabular}
\end{table}

The 3B $\rightarrow$ 7B step is \emph{flat}; the 7B $\rightarrow$ 8B step is a
generation change rather than a scale change. More informative than the aggregate:
\textbf{five cases fail on all three models}. The score differences arise from
other cases fluctuating, while this core remains invariant.

The value of this measurement is that it separates \emph{architectural limits} ---
addressable by design --- from \emph{capability limits} --- requiring a different
model. The cases currently blocking the system fall into neither category
cleanly, but they are demonstrably not resolved by increasing parameter count.

\subsection{Finding 4 --- Quantitative evidence of overfitting in rule design}
\label{sec:finding-overfit}

During development we introduced two rules to handle observed failure cases. Both
were written \emph{generically}, containing no entity names and no
language-specific patterns, and therefore satisfied our no-hard-coding constraint
in form.

We subsequently wrote \emph{falsifiers} --- fixtures designed specifically to
refute the specifications of those two rules. For instance, \emph{``Thủ đô của
nước đó nằm ở đâu?''} (``Where is the capital of that country?'') is a
location question whose jurisdiction is an unresolved reference, so the correct
action is \ACT{ask}. Both rules suppressed this legitimate clarification.

Removing the two rules produced the effect shown in Table~\ref{tab:overfit}.

\begin{table}[t]
\caption{Effect of removing the two unvalidated rules}
\label{tab:overfit}
\centering
\footnotesize
\begin{tabular}{@{}lr@{}}
\toprule
Evaluation suite & Effect \\
\midrule
Tuned    & $-1$ case \\
Held-out & unchanged \\
Safety   & unchanged \\
\bottomrule
\end{tabular}
\end{table}

The entire value of both rules resided on the very dataset against which they had
been designed, and was zero on unseen data.

\emph{Methodological implication.} A no-hard-coding constraint is \textbf{necessary
but not sufficient}. A generically phrased rule can still be a fixture fit if it
was \emph{selected} by observing which cases it repairs. Sufficiency requires a
held-out suite together with a mandatory falsifier for each rule.

%==============================================================================
\section{Limitations}
\label{sec:limitations}

We state the shortfalls explicitly.

\begin{enumerate}
  \item \textbf{Two wrong answers remain} on the independent test suites: an
        inverted subject--object question is accepted (the case in
        Section~\ref{sec:motivating}), and a question with an unresolved time
        reference is answered. Both carry valid citations and verbatim quotes,
        meaning they pass the entire validation chain of
        Section~\ref{sec:code-checks}. The system \emph{does not meet its design
        objective}.
  \item \textbf{The relation-reversal leak rate is approximately $30\%$ and is
        intermittent}: the same entity pair, differently phrased, is sometimes
        blocked and sometimes not.
  \item \textbf{Cost increases by roughly $1.9\times$} in model calls and
        $2.1\times$ in median latency.
  \item \textbf{Recall falls by $30\%$ relative} to the baseline; approximately
        half is attributable to the retriever ceiling and half to strict
        validation (Section~\ref{sec:recall-decomp}).
  \item \textbf{Results are computed on a 200/2000 sample}, with an approximate
        margin of $\pm 3.5$ percentage points. A full 2{,}000-question run
        requires roughly 20 hours on local hardware.
  \item \texttt{answer\_f1} \textbf{is lower than Prompt-Safe RAG} ($0.126$ versus
        $0.187$). This partly reflects lower recall and partly a metric artefact:
        PonyGuard emits terse verbatim quotes, which are penalised by token-level
        F1 against freely phrased gold answers.
  \item \textbf{The cross-provider capability comparison is incomplete} due to API
        quota limits, so we cannot yet report whether a stronger model closes the
        two remaining failures.
\end{enumerate}

%==============================================================================
\section{Future Work}

The intermittency of the relation-reversal leak
(Section~\ref{sec:limitations}, item~2) suggests a direct remedy:
\textbf{sample the judgement $k$ times and require unanimity}. If a failure occurs
in $30\%$ of trials, requiring five-way agreement drives the leak rate down
sharply, and the cost is incurred only on the branch about to answer.

Three further directions follow from the analysis. First, converting the guards
from fail-open to \textbf{fail-closed}, treated as a single measurable policy
parameter rather than as scattered defaults, and reporting the resulting
precision--recall curve. Second, \textbf{verifying from the answer side} rather
than the source side: composing a complete proposition from the produced answer
and asking whether the source entails it, since adjudicating one finished
proposition is an easier task than extracting semantic roles from a passage ---
an approach that two source-side attempts failed to make work. Third,
\textbf{two-stage retrieval} (broad retrieval followed by reranking), motivated by
the observation that the retriever ceiling of approximately $52\%$ upper-bounds
every system in the comparison.

%==============================================================================
\section{Conclusion}

We presented PonyGuard-ViQA, a Vietnamese question-answering architecture that
relocates the answer-acceptance decision from the language model into
deterministic validation code. On UIT-ViQuAD 2.0 the system reduces the rate of
answering questions that should be refused from $0.402$ to $0.134$ while raising
conditional precision from $0.530$ to $0.705$, indicating that the improvement
derives from better discrimination rather than from increased conservatism. The
cost is reduced recall and approximately $1.9\times$ compute.

Beyond the quantitative result, we regard the methodological findings as the more
transferable contribution: that verifying a derived relation with a
direct-verification criterion is ill-posed; that a specific failure class does not
diminish with model scale; and that a generically phrased rule can still
constitute a fixture fit if it was selected by observing which cases it repairs.

The system does not meet its design objective --- two wrong answers with valid
citations remain on the independent test suites. We report this together with the
localised cause and the intended remedy, rather than restricting the report to
favourable results.

%==============================================================================
\section*{Author Contributions}

This report was produced by \textbf{Group 3}. All five members participated in
the weekly design reviews, in the failure-analysis sessions from which the four
findings in Section~VII emerged, and in the writing of this report. Primary
areas of responsibility were divided as follows.

\begin{table}[h]
\caption{Division of responsibility}
\label{tab:contrib}
\centering
\footnotesize
\begin{tabular}{@{}lp{4.7cm}@{}}
\toprule
Member & Primary responsibility \\
\midrule
\textbf{Ngo Anh Duc}
  & Problem formulation and motivation; the three-way action taxonomy of
    Section~II; asymmetric error-cost analysis; curation of the motivating
    failure case. \\[2pt]
\textbf{Pham Ngoc Hoa}
  & Experimental design; the controlled three-system comparison; dataset
    preparation, frozen-split hashing, and retrieval configuration; the
    no-hard-coding constraint and its enforcement. \\[2pt]
\textbf{Trinh Duc Duong}
  & Pipeline architecture and implementation; the intent contract and its
    deterministic span validation; the seven-step validation chain and
    answer-plan locking; the inference proof path and direction check. \\[2pt]
\textbf{Do Quang Hiep}
  & Evaluation methodology; design of the tuned / held-out / safety suites and
    their positive controls; falsifier construction; the category-error and
    overfitting findings (Sections~VII-A and~VII-D). \\[2pt]
\textbf{To Thanh Hai}
  & Benchmark execution and metric implementation; the discrimination analysis
    of Section~VI-C and the recall decomposition of Section~VI-D; the
    model-scale study (Section~VII-C); limitations and future work. \\
\bottomrule
\end{tabular}
\end{table}

We note that the boundaries above are approximate. Several results reported here
--- in particular the two remaining safety failures documented in
Section~VIII --- were identified collaboratively during joint review rather than
by any single member.

%==============================================================================
\section*{Reproducibility}

All source code, versioned prompt contracts, configuration, and test suites are
contained in the project repository. Each run records the prompt versions, the
pinned model, and the provider that actually served every call, so results are
traceable. The test suite comprises 136 deterministic tests using scripted models,
which validate contracts rather than model competence.

To reproduce the main results:

\begin{lstlisting}[style=contract]
make test && make behavior-test
python scripts/evaluate_all.py --limit 200 --model local
\end{lstlisting}

%==============================================================================
\begin{thebibliography}{00}

\bibitem{viquad} K.~V. Nguyen \emph{et al.}, ``A Vietnamese dataset for
evaluating machine reading comprehension,'' in \emph{Proc. COLING}, 2020.

\bibitem{e5} L.~Wang \emph{et al.}, ``Multilingual E5 text embeddings,''
\emph{arXiv preprint arXiv:2402.05672}, 2024.

\bibitem{faiss} J.~Johnson, M.~Douze, and H.~J\'egou, ``Billion-scale similarity
search with GPUs,'' \emph{IEEE Trans. Big Data}, vol.~7, no.~3, pp.~535--547,
2021.

\bibitem{rag} P.~Lewis \emph{et al.}, ``Retrieval-augmented generation for
knowledge-intensive NLP tasks,'' in \emph{Proc. NeurIPS}, 2020.

\bibitem{selective} Y.~Geifman and R.~El-Yaniv, ``Selective classification for
deep neural networks,'' in \emph{Proc. NeurIPS}, 2017.

\bibitem{attribution} B.~Bohnet \emph{et al.}, ``Attributed question answering:
Evaluation and modeling for attributed large language models,'' \emph{arXiv
preprint arXiv:2212.08037}, 2022.

\bibitem{ambigqa} S.~Min, J.~Michael, H.~Hajishirzi, and L.~Zettlemoyer,
``AmbigQA: Answering ambiguous open-domain questions,'' in \emph{Proc. EMNLP},
2020.

\bibitem{qwen} Qwen Team, ``Qwen3 technical report,'' \emph{arXiv preprint},
2025.

\end{thebibliography}

\end{document}
